\section{Introduction}
\label{sec:intro}

AI agents that can browse the web on a user's behalf are rapidly
moving from research into real-world deployment
\citep{OpenAI2025Operator,anthropic2024computeruse,tamkin2024clioprivacypreservinginsightsrealworld,andreux2025surferhmeetsholo1costefficient,zhang2025surveylargelanguagemodel,nguyen-etal-2025-gui}.
To complete tasks like applying for jobs, purchasing goods, or paying
bills, these agents are given the user's full personally identifiable
information (PII): legal name, home address, payment card details,
government identifiers, account credentials, and API keys. The agent
stores this profile and uses it whenever a form needs to be filled or
a credential needs to be submitted.

This creates a serious security risk. The agent has broad access to
the web and, at the same time, holds sensitive personal data in its
context. Any malicious website it visits can potentially trick the
agent into handing that data over. This is a web-native version of
the confused-deputy problem
\citep{Hardy1988TheCD,Ferrag2025FromPI,Triedman2025MultiAgentSE,fu2024impromptertrickingllmagents,kim2025promptflowintegrityprevent,ji2026tamingvariousprivilegeescalation}. The agent cannot tell the difference between instructions from the user and instructions embedded in a malicious web page, and \textit{social-engineering attacks} exploit exactly this confusion. The web already hosts large amounts of phishing and manipulation infrastructure built to extract PII from human users~\citep{Hamadouche_Boudraa_Gasmi_2024,10.1145/3359183,10.1145/3173574.3174108,deceptivedesignpatterns,Bsch2016TalesFT}.
A human can pause, re-read, or close the tab. An agent is built to
finish its task and will keep filling forms and submitting data
without stopping to question whether it should. Capability benchmarks
confirm that frontier agents are now reliable enough to complete
multi-step web tasks end-to-end: \citet{garg2026real} report 41\%
success on complex state-changing tasks, and \textit{the same capability that
makes agents useful also makes them a predictable target for data
theft.}

\textbf{The measurement gap.} There is a growing body of work on the robustness of web agents
against adversarial attacks, but existing benchmarks measure the
wrong thing. They track whether the agent clicked a malicious link
or deviated from an expected trajectory, rather than whether the
agent actually submitted sensitive data to an attacker
(Table~\ref{tab:related-matrix}, \S\ref{sec:related})
\citep{cuvin2026decepticon,ersoy2026trickyarena,chen2026obviousinvisiblethreat,debenedetti2024agentdogo,korgul2025trap,zharmagambetov2026agentdam,xinyi2026webtrappark}.
What matters in practice is whether sensitive fields were submitted
to an attacker-controlled endpoint: that is the point of no return,
and it is the only outcome that the agent's own refusal can prevent.
To our knowledge, no existing benchmark combines (i)~attacker-controlled
websites that actively try to trick the agent into giving up PII,
(ii)~a realistic PII profile covering critical, high, and medium
sensitivity data, (iii)~field-level data submission as the primary
outcome measure, and (iv)~controlled environment pairs that let us
identify which attack design factors actually drive leakage.

\textbf{This work.} We introduce \textsc{Scammer4U}, a benchmark of 91 hand-crafted
attacker-controlled websites plus 10 benign baseline versions, built
to answer one question: \emph{when a web agent encounters a
social-engineering attack, will it give up the user's PII?} The
environments cover eight attack strategies across 16 website
categories, organized on an 8-axis taxonomy mapped to published
threat catalogues
\citep{mitre_attack_citation,owasp_llm_citation,enisa_citation}.
60 of the 91 environments are paired siblings that differ on exactly
one axis, so any difference in leakage can be attributed to that
single factor. We test four frontier model families (GPT-5 mini,
Claude Haiku 4.5, Gemini 3 Flash, Llama 4 Scout) under four
mitigation conditions ranging from no guidance to pre-submission
reflection, with the full analysis plan pre-registered before any
data was collected (\texttt{prereg-v2-start}).

Our experiments reveal three findings. Frontier agents leak
critical-tier PII at high rates even under prompt-level mitigations,
and the degree of protection varies sharply across model families.
More importantly, we find a detection--action gap: agents that
explicitly flag a site as suspicious in their own reasoning still
frequently go on to submit critical PII. This tells us that asking
agents to reason about trust before submitting is not a reliable
defense on its own. Effective protection will need to operate
outside the agent's reasoning loop entirely. We release
\textsc{Scammer4U} with a pre-registered analysis plan and complete
run logs so that future work can evaluate new models without
rebuilding the benchmark.